% ══════════════════════════════════════════════════════════════════
% CHAPITRE I — Présentation de l'organisme d'accueil
% ══════════════════════════════════════════════════════════════════
\newpage
\setcounter{chapter}{0}
\chapter{Présentation de l'organisme d'accueil}
\vspace{1cm}
\minitoc
\newpage

\textbf{Introduction}

Dans le cadre de l'industrie financière mondiale, le traitement sécurisé et efficace des transactions est d'une importance capitale. HPS Group s'inscrit dans ce contexte en tant qu'acteur majeur fournissant des solutions de paiement électronique de pointe. Ce chapitre a pour objectif de présenter la société d'accueil, HPS, ainsi que sa filiale ACPQUALIFE. Nous examinerons la structure organisationnelle, les secteurs d'activité, et les valeurs fondamentales de l'entreprise, afin de comprendre l'environnement professionnel dans lequel s'est déroulé ce stage de fin d'études et les exigences de qualité rigoureuses qui y sont appliquées.
\vspace{10pt}

\section{HPS Group}

HPS (Hightech Payment Systems) est un groupe compté parmi les leaders mondiaux dans l'édition et le développement des solutions de paiements électroniques, qui sont dédiées aux institutions financières. Son siège se trouve à Nearshore Park de Casablanca au Maroc HPS fournit des solutions progicielles basées sur la famille de produits PowerCARD, ainsi que les différents services associés, tels que l'implémentation du système et sa maintenance. Elle conçoit et fournit des solutions complètes, modulaires et intégrées. Grâce à PowerCARD, hautement paramétrable, HPS traite tous les types de cartes (cartes de crédit, cartes de débit, cartes prépayées, cartes de fidélité, cartes d'entreprises, ...), via tous les canaux (GAB, TPE, internet et mobile, ...) et pour toutes les catégories de commerçants (banques, sociétés, compagnies, opérateurs, ...)

\section{Historique d'HPS}

Fondée en 1995 par un groupe de consultants et d'experts, HPS est un fournisseur de solution monétique permettant la gestion du paiement électronique multi canal prenant en charge le terminal de paiement électronique, le guichet automatique bancaire, Internet et le mobile. Elle est certifiée par les principaux organismes internationaux tels que Visa, MasterCard, Diners Club, JCB et American Express.

HPS a progressivement étendu sa présence géographique en créant HPS Dubaï en 2003, puis HPS Europe en 2008. L'objectif premier de ces deux bureaux était de développer l'activité commerciale dans ces régions en délocalisant une force de vente dédiée, chargée de la commercialisation en direct et de l'animation d'un réseau de partenaires régionaux. En 2017, deux nouveaux bureaux sont ouverts à Singapore et Lausanne.

\section{Secteur d'activité}

Le progiciel PowerCARD offert par HPS est une solution de paiement électronique évoluée et hautement sécurisée. Il toute la chaîne de valeur du paiement et offre aux clients une panoplie de solutions à forte valeur ajoutée pour piloter et gérer leurs risques :

\begin{itemize}[leftmargin=*]
\item \textbf{PowerCARD-eSecure} : étape supplémentaire pour les paiements en ligne afin de lier une autorisation financière avec une authentification en ligne (ACS - 3D Secure).
\item \textbf{PowerCARD-ACH} : solution pour la gestion des autorisations et compensation entre différentes chambres.
\item \textbf{PowerCARD-Issuer} : services complets pour l'émission et la gestion de toutes les cartes sous tous les formats.
\item \textbf{PowerCARD-ATM} : solution self-service globale qui permet aux institutions financières et aux détaillants une gestion et une optimisation des GAB.
\item \textbf{PowerCARD-xPOS} : système de gestion des demandes d'autorisation et de gestion des transactions pour tout type de terminal de paiement électronique.
\item \textbf{PowerCARD-Tokenisation} : solution complète qui supporte l'émission, le provisionnement et le stockage de jetons de la part d'un demandeur de jeton.
\item \textbf{PowerCARD-Fraud} : solution de gestion de fraude autonome, qui contrôle en temps réel ou différé les autorisations venant d'un quelconque système.
\item \textbf{PowerCARD-eCommerce} : solution multi-commerçants et multi-acquéreurs pour une gestion de paiement e-commerce.
\item \textbf{PowerCARD-Switch} : routing, stand-in et autorisation dans un environnement à haute disponibilité.
\item \textbf{PowerCARD-Dashboard} : ensemble de tableaux de bord qui fournissent une vision claire aux utilisateurs en se basant sur leurs indicateurs clés de performance.
\item \textbf{PowerCARD-Acquirer} : plateforme globale de gestion des commerçants permettant aux acquéreurs d'adapter des solutions spécifiques.
\item \textbf{PowerCARD-WebPublisher} : solution portail web avec des fonctions administratives efficaces, qui permettent de gérer et de personnaliser de nombreux portails web.
\item \textbf{Retailer Open Payment Platform} : plateforme offrant un parcours client sans faille et sécurisé, acceptant les transactions de tout canal de paiement.
\end{itemize}

\section{Stratégie}

HPS a défini 4 axes stratégiques pour porter sa vision et lui permettre de maintenir sa position d'entreprise mondialement reconnue dans l'industrie du paiement électronique : privilégier une croissance durable, développer des solutions innovantes, renforcer l'excellence opérationnelle et impliquer fortement les collaborateurs.

\section{Mission de l'entreprise}

Après plus de 20 ans d'activité, HPS est l'un des leaders mondiaux de l'industrie. Il accompagne plus de 500 institutions à travers 98 pays dans le monde en développant des offres à même de couvrir tous les aspects liés à l'activité de paiement multicanal. Ces offres sont structurées autour de 3 Business Units qui agissent de manière complémentaire : HPS Solutions, HPS Processing et HPS Services. La figure suivante représente l'organigramme d'HPS.

\begin{figure}[H]
\centering
\includegraphics[width=0.6\textwidth]{page12_img02.jpeg}
\caption{Organigramme d'HPS}
\end{figure}

\subsection{HPS Solutions}

Entité historique du Groupe, HPS Solutions conçoit, développe, commercialise, installe et supporte une suite de solutions de paiement électronique multicanal. Leur suite logicielle PowerCARD propose des solutions complètes, paramétrables et évolutives, à même de répondre aux besoins les plus spécifiques des utilisateurs. Évalué « Positive » par le Gartner durant 3 années consécutives, PowerCARD se positionne dans le Top 5 mondial des meilleures solutions de l'industrie monétique.

\subsection{HPS Services}

Au sein du Groupe, les activités de services IT sont portées par HPS Services, qui \oe uvre pour la maîtrise et la performance des systèmes d'informations pour le compte de grands groupes internationaux. HPS Services repose sur la structure et les équipes de la société Acpqualife, créée en 2002 et filiale à 100\% du Groupe HPS depuis 2010.

Les activités de HPS Services se basent sur l'expertise internationale de plus de 75 consultants dans les domaines du test et de la qualification logicielle. HPS Services offre également des services d'assistance à maîtrise d'ouvrage dans le cadre de projets de transformation des SI et des services d'ingénierie IT autour des technologies JEE, Spring, JBOSS, Angular Js ou encore JSF, et associés à l'implémentation de méthodologies de type agile.

\subsection{HPS Processing}

HPS Processing, créée en 2016, met en \oe uvre des moyens humains et matériels, des services à forte valeur ajoutée et son expertise autour de PowerCARD pour offrir à ses clients les solutions du Groupe en mode externalisé. L'activité Processing couvre 3 secteurs d'activité distincts :

\begin{itemize}[leftmargin=*]
\item Le Processing de solutions de paiement qui s'adresse aux institutions financières et opérateurs de paiement pour gérer l'ensemble de leurs activités monétiques.
\item La gestion des activités de Switching à l'échelle nationale ou régionale qui consiste à gérer les demandes d'autorisation et à traiter la compensation des flux entre différents opérateurs.
\item Le Processing de solutions pour les acteurs de micro finance pour la gestion complète de leur activité et la dématérialisation des flux financiers liés au déblocage et au remboursement des crédits.
\end{itemize}

\section{Chiffres clés}

Avec une croissance continue et de récentes acquisitions, HPS Group compte plus de 1 500 collaborateurs et dessert plus de 500 institutions dans plus de 98 pays sur 5 continents.

\begin{table}[H]
\centering
\caption{Chiffres clés d'HPS (2024)}
\begin{tabularx}{0.9\textwidth}{l X}
\toprule
\textbf{Indicateur} & \textbf{Valeur} \\
\midrule
Clients & Plus de 500 \\
Pays & Plus de 98 \\
Continents & 5 \\
Employés & Plus de 1 500 \\
Bureaux & Casablanca, Paris, Aix-en-Provence, Dublin, Dubaï, Jordanie, Singapour, Pune, Sydney, Le Cap, Maurice, Montréal \\
\bottomrule
\end{tabularx}
\end{table}

\section{ACPQUALIFE - HPS Group}

ACPQUALIFE est une société spécialisée dans les systèmes d'informations (SI), particulièrement dans le domaine de test et qualification logicielle, relatif à la monétique et à la gestion des transactions électroniques. ACPQUALIFE met toute son expertise pour des systèmes d'informations performants. Elle s'illustre notamment pour son expertise internationale dans le domaine du Test et de la Qualification logicielle. ACPQUALIFE est un expert des systèmes d'informations, il fournit une offre de services complète et sur-mesure, possède une structure entièrement dédiée aux métiers du test et un centre d'expertise dans le secteur des moyens de paiement électroniques. Elle dispense une formation pragmatique et certifiante.

\section{Historique}

La société ACPQUALIFE a été créée en 2002, suite à la fusion de deux entreprises :
\begin{itemize}[leftmargin=*]
\item ACP (Architect \& Consulting Partners), fondée par M. Christian COURTES, son Directeur Général actuel, en charge des activités services,
\item QUALIFE, spécialisée dans la qualification de systèmes d'information, fondée par M. Marc DURUPT, son Directeur Associé actuel, en charge des expertises sur les métiers du test, en collaboration avec Mme Véronique THEAULT, sa Directrice Associée actuelle, en charge des offres de qualification.
\end{itemize}

En 2010, la société ACPQUALIFE a fusionnée avec le groupe HPS (Hightech Payment Systems), spécialisé dans les paiements électroniques, dont le siège se trouve à Casablanca au Maroc.

\section{Mission}

La société ACPQUALIFE a pour mission principale d'apporter à ses clients des solutions performantes et pérennes ; afin de maîtriser leur système d'information. Et ce, depuis la conception jusqu'à la mise en production, en passant par les phases d'architectures, de réalisations et de tests.

ACPQUALIFE veille sur la sécurité des transactions électroniques, l'adaptabilité simplifiée aux standards internationaux, le développement des solutions mobiles, \ldots

\section{Valeurs (ACPQ)}

ACPQUALIFE fonde son identité et affirme son image auprès de ses clients et collaborateurs à travers quatre valeurs principales :
\begin{itemize}[leftmargin=*]
\item \textbf{Anticipation} : une capacité d'innovation qui permet de proposer aux clients des solutions adaptées et performantes,
\item \textbf{Professionnalisme} : des propositions bâties sur des savoir-faire éprouvés qui sont le résultat d'une capitalisation des expériences,
\item \textbf{Contrôle} : Un développement ancré autour de la culture de l'engagement qui permet une optimisation de la maîtrise des interventions,
\item \textbf{Qualité} : Dans toutes les interventions, application de la démarche « Qualife », qui vise à garantir un véritable label de qualité.
\end{itemize}

\section{Expertises}

Propose une offre de services complète et sur-mesure dont la maîtrise technique est mise au service des activités métiers de ses clients,
\begin{itemize}[leftmargin=*]
\item Maîtrise des réalisations dans le domaine des NTIC et de l'Ingénierie Mobile (Architecture, Méthodologie, Réalisation, Recette)
\item Maîtrise des environnements de production et des processus à mettre en \oe uvre pour garantir leur productivité (Expertise ITIL, Mise en \OE uvre de SLA, Devops, Cloud, Hadoop\ldots)
\item Expertise dans le domaine des Tests et de la Qualification Logicielle : Conseil, Formation, Mise en \OE uvre, Assistance MOA 2 et MOE3, Expertise Outils, Tierce Recette Applicative.
\end{itemize}

\section{Services}

\begin{itemize}[leftmargin=*]
\item \textbf{Maîtrise d'ouvrage AMOA} : ACPQUALIFE est un partenaire qui apporte des idées innovantes dans le cadre des évolutions et de la transformation des SI, en s'appuyant sur une parfaite connaissance du secteur métier ainsi que sur son expertise technique et méthodologique,
\item \textbf{Qualification} : ACPQUALIFE possède un département entièrement dédié aux métiers du test. Il s'appuie notamment sur « Qualitest », une méthodologie pragmatique, rodée et adaptable,
\item \textbf{Ingénierie d'Infrastructure} : ACPQUALIFE accompagne ses clients dans la continuité de services et la performance de leur Système d'Information,
\item \textbf{Ingénierie IT} : Au travers d'une activité concentrée sur les Nouvelles Technologies de l'Information et des Communications, ACPQUALIFE a acquis une spécialisation reconnue dans la réalisation d'applications en Architectures.
\end{itemize}

\section{Chiffres Clés}

\begin{table}[H]
\centering
\caption{Chiffres clés ACPQUALIFE - HPS Group}
\begin{tabularx}{0.9\textwidth}{l X}
\toprule
\textbf{ACPQUALIFE - HPS Group} & \\
\midrule
Adresse & ACPQUALIFE, 805 Av. J.R.G Gautier De La Lauzière 13290 Aix En Provence - France \\
Activité & Gestion d'installations informatiques \\
Forme juridique & Société par actions simplifiée \\
Date immatriculation & Le 09-07-2002 \\
Tranche d'effectif & 50 à 99 salariés \\
Chiffre d'affaires 2016 & 12 938 700,00 \euro \\
Implantation & Aix en Provence, Paris, Toulouse, Lausanne, Lille et Sophia Antipolis. \\
\midrule
\multicolumn{2}{l}{Plus de 100 consultants} \\
\multicolumn{2}{l}{Plus 28\% de croissance} \\
\multicolumn{2}{l}{Plus de 25\% d'activité à l'international} \\
\bottomrule
\end{tabularx}
\end{table}

\vspace{10pt}
\textbf{Conclusion}

En somme, HPS et sa filiale ACPQUALIFE constituent un environnement hautement spécialisé et exigeant, dédié à la performance et à la sécurité des systèmes d'information monétiques. L'expertise internationale du groupe et sa culture de la qualité fournissent un cadre idéal pour l'apprentissage et l'application des meilleures pratiques en ingénierie logicielle. La compréhension de cette organisation permet de mieux appréhender les défis techniques et opérationnels du projet d'automatisation des tests, dont le contexte et les objectifs spécifiques seront détaillés dans le chapitre suivant.
